\uuid{38HI}
\exo7id{7031}
\titre{exo7 7031}
\auteur{megy}
\organisation{exo7}
\datecreate{2016-10-28}
\isIndication{false}
\isCorrection{true}
\chapitre{Déterminant, système linéaire}
\sousChapitre{Système linéaire, rang}
\cpge{mpsi-mp2i.calcul-algebrique-trigonometrie.b, pcsi.calcul-algebrique-trigonometrie.b, ptsi.calcul-algebrique-trigonometrie.b, mpsi-mp2i.calcul-matriciel-systemes.c, pcsi.calcul-matriciel-systemes-lineaires.c, ptsi.calcul-matriciel-systemes-lineaires.c, bcpst1.nombres-reels, bcpst1.systemes-lineaires.a}
\module{Algèbre}
\niveau{L2}
\difficulte{}

\contenu{
\question{
Résoudre dans $\R_+^*$ le système suivant:
\[
\begin{cases}
xyz=1\\
xy^2z^4=2\\
xy^3z^9=3
\end{cases}
\]
}
\reponse{
On trouve une unique solution $\left(\frac{3}{8},\frac{16}{3\sqrt 3},\frac{1}{\sqrt 3}\right)$.
}
}
